\documentclass[10pt,a4paper]{amsart}
\usepackage[margin=19mm]{geometry}
\usepackage{amsmath,amssymb,mathtools}
\usepackage[scr=rsfs,cal=euler]{mathalfa}
\usepackage{xfrac}
\usepackage[unicode,hidelinks,bookmarks=false]{hyperref}
\newcommand{\ie}{i.e.\ }
\newcommand{\cf}{cf.\ }
\newcommand{\resp}{respectively\ }
\newcommand{\Subsection}{\subsection}
\providecommand{\nobreakdash}{\nobreak\hbox{-}}
\setlength{\emergencystretch}{2em}
\multlinegap0pt
\usepackage[shortlabels]{enumitem}
\usepackage{amssymb}
\usepackage{mathtools}
\usepackage{xcolor}
\newtheorem{definition}{Definition}
\newtheorem{lemma}{Lemma}
\newtheorem{proposition}{Proposition}
\newcommand{\Ex}{\mathcal E}
\newcommand{\R}{\mathbb R}
\DeclareMathOperator{\Ric}{Ric}
\DeclareMathOperator{\Vol}{Vol}
\newcommand{\eps}{\varepsilon}
\title{Universal Volume Growth Bounds from Positive\\ Intermediate Curvature}
\author{Gioacchino Antonelli}
\date{Source: arXiv:2608.14507v2 (CC BY 4.0)}
\begin{document}
\maketitle

\noindent\emph{Source and licence.} Universal Volume Growth Bounds from Positive\\ Intermediate Curvature. Version arXiv:2608.14507v2 (CC BY 4.0), distributed by arXiv under the Creative Commons Attribution 4.0 International licence, http://creativecommons.org/licenses/by/4.0/, as recorded in the arXiv licence metadata for this exact version and shown on the abstract page https://arxiv.org/abs/2608.14507v2. Full licence notice: \texttt{licenses/arxiv-2608.14507-CC-BY-4.0.txt}.\par\smallskip

\noindent\textit{Editorial scope.} Counted unit: the article's proposition prop:rank-improvement (source0.tex lines 792--804). The article's complete original proof of it is delivered with it (source0.tex lines 806--972, one \texttt{proof} environment of 167 lines). No mathematics is written, repaired or re-derived: the statement and the proof are the article's own lines, in the article's own notation, and no retained line is substituted or re-flowed.  Retained uncounted as the source-local context the proof needs: lines 344--353; lines 554--563; lines 611--631.  Nothing was omitted from any retained span, so the package carries no omission record.  No retained line contains a citation, so the wrapper carries no bibliography and no bibliography entry.  No retained line contains a figure environment, an image-inclusion command or drawing code, so no image is delivered and the package declares no asset dependency.  Editorial formatting only, in addition: an AMS \texttt{amsart} preamble that restates the article's own declarations and macros used by the retained lines, namely the environments \texttt{definition}, \texttt{lemma}, \texttt{proposition} on separate counters, together with the standard packages those lines need.  The retained environments print in this document as Definition 1, Lemma 1, Proposition 1 in order of appearance, whereas the article prints the same items with the numbering of its own full document.  The complete native source of the article as distributed in the arXiv e-print for this exact version is bundled unchanged as \texttt{sources/207674/source0.tex}.
\begin{definition}\label{def:metric-split}
Let $0<\eps<1$ and let $k,n$ be integers such that $0\leq k\leq n$.  A ball $B_r(x)$ is
\emph{$(k,\eps)$-split} if there are a pointed proper length space
$(Z,z)$ and a pointed $\eps r$-Gromov--Hausdorff approximation
\[
 B_{\eps^{-1}r}(x)\longrightarrow
 B_{\eps^{-1}r}((0,z))\subset\R^k\times Z.
\]
Note that every ball is $(0,\eps)$-split.
\end{definition}

\begin{lemma}\label{lem:finite-segment}
For every integer $n\geq 1$, every integer $0\leq k<n$, and $0<\eps<1$, there are
$L=L(n,k,\eps)>\eps^{-1}$ and
$0<\delta_s=\delta_s(n,k,\eps)<1$ with the following property.
Let $(M^n,g)$ be a complete Riemannian manifold, $s>0$, assume $\Ric\geq -\delta_ss^{-2}$, and let $(Z,d_Z,z_0)$ be a pointed proper geodesic space. Assume the closed ball $\overline B_{Ls}(q)$ is
$\delta_s s$-close to the closed $Ls$-ball centered at $(0,z_0)$ in
$\R^k\times Z$.  If $z_0$ is the midpoint of a minimizing 
segment $\gamma:[-Ls,Ls]\to Z$, then
$B_s(q)$ is $(k+1,\eps)$-split.
\end{lemma}

\begin{proposition}
\label{prop:critical-package}
Fix an integer $n\geq 1$, an integer $0\leq k<n$, and $0<A<\infty$, $0<\eta<1$.
There are $0<\delta_c=\delta_c(n,k,A,\eta)<1$,
$c=c(n)>0$, $C=C(k,n)<\infty$, and $c_0=c_0(n)>0$ with the following
property.  Suppose $(M^n,g,p)$ is complete, has $\Ric_g\geq -\delta_c$, and
\[
 d_{GH}\bigl(B_{\delta_c^{-1}}(p),
             B_{\delta_c^{-1}}^{\R^k}(0)\bigr)\leq\delta_c.
\]
Then there are a measurable set $G\subset B_{1/4}(p)$ and an integer $N\geq 1$ for which the following holds. For every $1\leq j\leq N$ there are $q_j\in B_{1/4}(p)$, $0<t_j\leq 1$, pairwise disjoint measurable $E_j\subset G$, pointed proper geodesic spaces $(Y_j,d_{Y_j},y_j)$, and $a_j\in Y_j$ such that:
\begin{enumerate}[label=\textup{(\roman*)},leftmargin=2.3em]
\item $G=\sqcup_{j=1}^N E_j$ and $\Vol(G)\geq c\Vol B_1(p)$;
\item $\sum_{j=1}^N t_j^k\leq C$, and, for every $1\leq j\leq N$, $E_j\subset B_{Ct_j}(q_j)$;
\item For every $1\leq j\leq N$, $d_{Y_j}(y_j,a_j)\geq c_0t_j$, and
\[
 d_{GH}\bigl(B_{At_j}(q_j),
             B_{At_j}^{\mathbb R^k\times Y_j}((0,y_j))\bigr)\leq\eta t_j.
\]
\end{enumerate}
\end{proposition}

\begin{proposition}
\label{prop:rank-improvement}
Fix an integer $n\geq2$, integers $m,k$ such that $0\leq m\leq n-2$, $0\leq k\leq m$, and fix
$0<\eps<1$.  There are
$0<\delta=\delta(n,m,k,\eps)<1$ and
$c=c(n,m,k,\eps)>0$ such that the following holds. Let $(M^n,g)$ be a complete Riemannian manifold, $p\in M$ and $r>0$. If
$\Ric\geq -\delta r^{-2}$ on $M$ and $B_r(p)$ is $(k,\delta)$-split, then there are
$q\in M$ and $0<s\leq r$ such that $B_s(q)$ is
$(k+1,\eps)$-split and
\begin{equation}\label{eq:rank-volume}
                        \Ex_m(q,s)\geq c\Ex_m(p,r).
\end{equation}
\end{proposition}

\begin{proof}
All constants denoted by $c,C,C'>0$ in this proof depend only on
$n,m,k,\varepsilon$ and may change from line to line. It is enough to prove the rescaled version with
$r=1$. So, from now on, let us assume $r=1$.
Notice
\[
                         \Ex_m(p,1)=\Vol B_1(p).
\]
Assume $B_{1}(p)$ is $(k,\delta)$-split. Let
\begin{equation}\label{eqn:Original}
 \Phi:B_{\delta^{-1}}(p)\longrightarrow
 B_{\delta^{-1}}((0,z))\subset\mathbb R^k\times Z
\end{equation}
be the $\delta$-approximation given by  Definition \ref{def:metric-split}.
\smallskip

The proof is divided into two cases. If the residual factor $Z$ has a long enough segment through $z$, Lemma \ref{lem:finite-segment} directly produces an additional Euclidean direction. Otherwise Proposition
\ref{prop:critical-package} gives a collection of
good cells $E_j$ in $M$ at different scales, with bounded diameter.  Since their scales are summable, we can select one cell carrying a
definite fraction of the $m$-volume ratio. A segment in the space $Y_j$ associated to such a cell supplies the additional Euclidean direction. Let us carefully perform this strategy.
\smallskip

Let $L(n,k,\varepsilon)>\varepsilon^{-1}>1,0<\delta_s(n,k,\varepsilon)<1$ be the constants in Lemma \ref{lem:finite-segment}, and let $c_0$ be the constant in Proposition \ref{prop:critical-package}. Choose
\[
 0<\vartheta\le c_0/(32L),\quad
 0<\eta\le\min\{1/2,\delta_s\vartheta/C_1\},\quad
 A>c_0/4+L\vartheta+2,\quad
 d_0=\delta_c/20,\quad R_0=4\delta_c^{-1},
\]
where $C_{1}$ is a large universal constant that controls restriction and recentering of a
Gromov--Hausdorff approximation, and $\delta_c(n,k,A,\eta)<1$ is the corresponding constant in
Proposition \ref{prop:critical-package}. Finally choose $\delta$ sufficiently
small so that $0<\delta\leq \min\{1,\delta_s,\delta_c\}$ and:
\begin{equation}\label{eqn:ChoiceOfDelta}
    1+\frac{3d_0}{4}\leq \delta^{-1}, \qquad C_1\delta\leq \frac{\delta_s d_0}{4L}, \qquad \delta_c^{-1}<\frac{\delta^{-1}+1}{10}, \quad C_1\delta\leq \frac{\delta_c}{2}.
\end{equation}

\smallskip
\noindent
\emph{Case 1: There is a long enough segment through $z$ in $Z$.}
Suppose that some $z'\in B_{R_0}(z)$ satisfies
$d_Z(z,z')\geq d_0$. Let $\gamma:[0,d_0]\to Z$ be the unit-speed initial segment of a minimizing geodesic from $z$ to $z'$. Set
\[
 \bar z:=\gamma(d_0/2),
 \qquad
 s_0:=\frac{d_0}{4L}.
\]
Since $Ls_0=d_0/4$, the restriction of $\gamma$ to
$[d_0/4,3d_0/4]$, reparametrized on $[-Ls_0,Ls_0]$, is a minimizing
segment centered at $\bar z$. Choose $q\in M$ corresponding to $(0,\bar z)$ under the approximation \eqref{eqn:Original}. First
$ \overline B_{Ls_0}((0,\bar z))
 \subset \overline B_{3d_0/4}((0,z))$,
so this ball lies inside the original comparison region since $1+3d_0/4\le \delta^{-1}$ (see \eqref{eqn:ChoiceOfDelta}). Restriction
and recentering give
\[
 d_{GH}\!\left(
 \overline B_{Ls_0}(q),
 \overline B_{Ls_0}((0,\bar z))
 \right)
 \leq C_1\delta,
\]
for a universal constant $C_1>0$.
Since
 $C_{1}\delta\leq\delta_s s_0$ (see \eqref{eqn:ChoiceOfDelta}), and $\mathrm{Ric}\geq -\delta\geq -\delta_ss_0^{-2}$ (using $s_0\leq 1$) all the hypotheses of Lemma \ref{lem:finite-segment} are
satisfied, and consequently $B_{s_0}(q)$ is
$(k+1,\varepsilon)$-split.

Since $p$ and $q$ correspond respectively to $(0,z)$ and
$(0,\bar z)$ in the approximation given by \eqref{eqn:Original},
\[
 d(p,q)\leq d_Z(z,\bar z)+C_{\rm 1}\delta
          =\frac{d_0}{2}+C_{\rm 1}\delta
          \leq C(n,m,k,\varepsilon),
\] 
and hence $B_1(p)\subset B_{C}(q)$ for a possibly larger constant $C>\max\{1,s_0\}$.
Bishop--Gromov monotonicity then gives, with $c:=c(n,m,k,\varepsilon)$ possibly changing from line to line
\[
 \begin{split}
 \Ex_m(q,s_0)
 &=\frac{\Vol B_{s_0}(q)}{s_0^m}
 \geq \frac{c}{s_0^m}\Vol B_C(q)
 \geq c\,\Vol B_1(p)
 =c\Ex_m(p,1).
 \end{split}
\]
We thus proved that $B_{s_0}(q)$ is $(k+1,\varepsilon)$-split and $\Ex_m(q,s_0)\geq c\Ex_m(p,1)$. Hence, the proof is concluded in the first case. Notice that we have $s_0=d_0/4L\leq 1=r$, as desired.

\smallskip
\noindent
\emph{Case 2: The residual factor $Z$ is macroscopically small.}
Suppose now that every point of $B_{R_0}(z)$ lies within $d_0$ of
$z$. Since $\delta_c^{-1}<R_0$, projection onto the Euclidean factor
gives
\[
 d_{GH}\!\left(
 B_{\delta_c^{-1}}^{\mathbb R^k\times Z}((0,z)),
 B_{\delta_c^{-1}}^{\mathbb R^k}(0)\right)\leq2d_0.
\]
Hence, using \eqref{eqn:Original} and the third inequality in \eqref{eqn:ChoiceOfDelta}
\[
 d_{GH}\!\left(
 B_{\delta_c^{-1}}(p),
 B_{\delta_c^{-1}}^{\mathbb R^k}(0)\right)\leq\delta_c,
\]
where we also used $C_1\delta\leq \delta_c/2$, see \eqref{eqn:ChoiceOfDelta}. Moreover $\mathrm{Ric}\geq - \delta\geq -\delta_c$. 
Apply Proposition \ref{prop:critical-package}, and denote its output
by $G$ and $\{q_j,t_j,E_j,Y_j,y_j,a_j\}_{j=1}^N$.
Since $k\leq m$ and $0<t_j\leq1$, we have $
                    \sum_{j=1}^N t_j^m
                    \leq\sum_{j=1}^N t_j^k\leq C$. 
Consequently, for some $j$
\[
 \frac{\Vol(E_j)}{t_j^m}
 \geq\frac{\Vol(G)}{\sum_{\ell=1}^N t_\ell^m}
 \geq c\Vol B_1(p).
\]
Since $E_j\subset B_{Ct_j}(q_j)$, we have the following inequality where the constant $c$ might change from line to line
\begin{equation}\label{eq:heavy-cell}
 \Ex_m(q_j,Ct_j)
 =\frac{\Vol B_{Ct_j}(q_j)}{(Ct_j)^m}
 \geq c\frac{\Vol(E_j)}{t_j^m}
 \geq c\Vol B_1(p)=c\Ex_m(p,1).
\end{equation}

Let $\gamma:[0,c_0t_j/2]\to Y_j$ be the initial part of a minimizing
geodesic from $y_j$ to $a_j$, and put $
                       \bar y:=\gamma(c_0t_j/4)$, and $s:=\vartheta t_j$.
Thus $\bar y$ is the midpoint of a segment of half-length $c_0t_j/4$.
Since $Ls=L\vartheta t_j\leq c_0t_j/32$, the central subsegment of
half-length $Ls$ satisfies the segment hypothesis of
Lemma \ref{lem:finite-segment}.  Moreover, by the choice of $A$
\[
 \overline B_{Ls}((0,\bar y))
 \subset
 \overline B_{(c_0/4+L\vartheta)t_j}((0,y_j))
 \subset \overline{B}_{At_j}((0,y_j)).
\]
Choose $q\in M$ corresponding to $(0,\bar y)$ in the product
approximation given by Proposition \ref{prop:critical-package} at $q_j$.  Restriction and recentering give
\[
 d_{GH}\!\left(
 \overline B_{Ls}(q),
 \overline B_{Ls}((0,\bar y))
 \right)
 \leq C_1\eta t_j
 \leq\delta_s\vartheta t_j=\delta_ss.
\]
Since $\mathrm{Ric}\geq - \delta\geq -\delta_ss^{-2}$ (using $s=\vartheta t_j\leq 1$), Lemma \ref{lem:finite-segment} therefore shows that
$B_s(q)=B_{\vartheta t_j}(q)$ is $(k+1,\varepsilon)$-split.

The proof is now concluded as in Case 1: we have $d(q,q_j)\leq Ct_j$, so
$B_{Ct_j}(q_j)\subset B_{C't_j}(q)$. Bishop--Gromov monotonicity and
\eqref{eq:heavy-cell} yield, possibly by changing constants $c,C$ from line to line,
\[
 \begin{split}
 \Ex_m(q,\vartheta t_j)
 &=
   \frac{\Vol B_{\vartheta t_j}(q)}{(\vartheta t_j)^m}\\
 &\geq
   \frac{c}{(\vartheta t_j)^m}
   \Vol B_{C't_j}(q)\geq c\frac{\Vol B_{Ct_j}(q_j)}{t_j^m}
 \geq c\Ex_m(q_j,Ct_j)
 \geq c\Ex_m(p,1),
 \end{split}
\]
thus concluding the proof. Notice that we also have $\vartheta t_j\leq 1=r$, as desired.
\end{proof}

\end{document}
